\documentclass[journal]{IEEEtran}

% --- LENGUAJE Y CODIFICACIÓN ---
\usepackage[utf8]{inputenc}   % Permite tildes y ñ directamente (UTF-8 es el estándar)
\usepackage[T1]{fontenc}      % Codificación de fuente para que las tildes se vean bien al copiar texto del PDF
\usepackage[spanish, es-tabla]{babel} % Configura el documento en español ("Cuadro" cambia a "Tabla")

% --- MATEMÁTICAS Y SÍMBOLOS ---
\usepackage{amsmath, amssymb, amsfonts} % El triunvirato de poder para fórmulas matemáticas

% --- GRÁFICOS Y TABLAS ---
\usepackage{graphicx}  % Para insertar imágenes (capturas de pantalla de su HMI)
\usepackage{float}     % Permite usar [H] para forzar la posición de imágenes y tablas
\usepackage{booktabs}  % Tablas de calidad profesional (sin líneas verticales, estilo IEEE)
\usepackage{multirow}  % Celdas que ocupan varias filas

% --- CÓDIGO FUENTE (PAQUETE NUEVO) ---
\usepackage{listings}  % Esencial para mostrar código con resaltado de sintaxis
\usepackage{xcolor}    % Para dar color al código

% Configuración de colores para el código (Estilo Dark Mode / Profesional)
\definecolor{codegreen}{rgb}{0,0.6,0}
\definecolor{codegray}{rgb}{0.5,0.5,0.5}
\definecolor{codepurple}{rgb}{0.58,0,0.82}
\definecolor{backcolour}{rgb}{0.95,0.95,0.92}

\lstset{
    commentstyle=\color{codegreen},
    keywordstyle=\color{magenta},
    numberstyle=\tiny\color{codegray},
    stringstyle=\color{codepurple},
    basicstyle=\ttfamily\footnotesize,
    breakatwhitespace=false,         
    breaklines=true,                 
    captionpos=b,                    
    keepspaces=true,                 
    numbers=left,                    
    numbersep=5pt,                  
    showspaces=false,                
    showstringspaces=false,
    showtabs=false,                  
    tabsize=2,
    language=C++ % Por defecto para Arduino/ESP32
}

% --- HIPERVÍNCULOS (CONFIGURACIÓN SOBERANA) ---
\usepackage{hyperref}
\hypersetup{
    colorlinks=true,
    linkcolor=blue,
    filecolor=magenta,      
    urlcolor=cyan,
    pdftitle={Reporte de Actividad - [Su Nombre]},
}

% --- MEJORAS TIPOGRÁFICAS ---
\usepackage{microtype} % Optimiza el espaciado entre letras para evitar huecos blancos feos
\usepackage{csquotes}  % Manejo inteligente de comillas según el idioma

\begin{document}

% --- INICIO DE LA PORTADA (Rúbrica 1%) ---
\begin{titlepage}
    \centering
    \vspace*{0.5cm}
    
    % --- Escudo de la UV ---
    % Asegúrese de tener el archivo "escudo_uv.png" (o .jpg) en su carpeta del proyecto
    \includegraphics[width=0.3\textwidth]{UV.png} \\ 
    \vspace{0.8cm}
    
    {\huge \textbf{Universidad Veracruzana}} \\
    \vspace{0.4cm}
    {\Large Facultad de Instrumentación Electrónica} \\ % Corregido: Instrumentación
    \vspace{0.2cm}
    {\large Ingeniería Biomédica} \\
    
    \vfill
    
    % Título de la Experiencia Educativa y Actividad
    {\Large \textbf{EE: Interfaces Hombre-Máquina}} \\
    \vspace{0.8cm}
    {\Large \textbf{Actividad 01: Comunicación RS-232 \\ con Interfaz RX-TX Simple}} \\
    
    \vspace{1.5cm}
e   
    % Información de los Integrantes
    \begin{flushright} % Alineado a la derecha para un toque más profesional
        \textbf{Profesor:} \\
        Varela Lara Luis Julian \\
        \vspace{0.5cm}
        \textbf{Integrantes del equipo:} \\
        Alfonso Gamboa Rubén \\
        Flores Montero Edsel Yetlanezi \\
    \end{flushright}
    
    \vfill
    
    {\large \textbf{Fecha:} \today} \\
    \vspace{1cm}
    {\small Xalapa, Enríquez, Ver.} % O la sede correspondiente
\end{titlepage}
% --- FIN DE LA PORTADA ---

\newpage
\pagenumbering{arabic} % Reiniciar numeración si es necesario

% Aquí empieza el formato IEEE real
\title{Implementación de Control de LEDs Mediante Interfaz HMI Virtual y Protocolo RS-232 No Bloqueante}
\author{Su Nombre}
\maketitle

\begin{abstract}
Este reporte documenta el desarrollo de un firmware...
\end{abstract}

\section{Introducción}
...
\end{document}
